\documentclass[10pt,a4paper]{amsart}
\usepackage[margin=19mm]{geometry}
\usepackage{amsmath,amssymb,mathtools}
\usepackage[scr=rsfs,cal=euler]{mathalfa}
\usepackage{xfrac}
\usepackage[unicode,hidelinks,bookmarks=false]{hyperref}
\newcommand{\ie}{i.e.\ }
\newcommand{\cf}{cf.\ }
\newcommand{\resp}{respectively\ }
\newcommand{\Subsection}{\subsection}
\providecommand{\nobreakdash}{\nobreak\hbox{-}}
\setlength{\emergencystretch}{2em}
\multlinegap0pt
\usepackage{bm}
\usepackage{bbm}
\usepackage{dsfont}
\usepackage{xspace}
\usepackage{cancel}
\usepackage{mathtools}
\usepackage{amsthm}
\makeatletter
\@ifundefined{thm}{\newtheorem{thm}{Thm}}{}
\@ifundefined{ass}{\newtheorem{ass}[thm]{Assumption}}{}
\@ifundefined{lem}{\newtheorem{lem}[thm]{Lemma}}{}
\@ifundefined{prop}{\newtheorem{prop}[thm]{Proposition}}{}
\makeatother
\newcommand{\N}{\mathbb{N}}
\newcommand{\R}{\mathbb{R}}
\newcommand{\cA}{\mathcal{A}}
\newcommand{\cY}{\mathcal{Y}}
\renewcommand{\d}{{\rm d}}
\title[Almost sure convergence rates of adaptive increasingly rare ]{Almost sure convergence rates of adaptive increasingly rare Markov chain Monte Carlo}
\author{Julian Hofstadler, Krzysztof Latuszynski, Gareth O. Roberts, Daniel Rudolf}
\date{Source: arXiv:2402.12122v2 (CC BY 4.0)}
\begin{document}
\maketitle

\noindent\emph{Source and licence.} Almost sure convergence rates of adaptive increasingly rare Markov chain Monte Carlo. Version arXiv:2402.12122v2 (CC BY 4.0), distributed under the Creative Commons Attribution 4.0 International licence, as recorded in the arXiv licence metadata for this exact version and shown on the abstract page https://arxiv.org/abs/2402.12122v2. Full rights notice: licenses/arxiv-2402.12122-CC-BY-4.0.txt.\par\smallskip

\noindent\textit{Editorial scope.} Counted unit: the Prop whose source label is \texttt{prop:poisson\_equation\_wasserstein\_setting} in the article (source0.tex lines 424--435, source label \texttt{prop:poisson\_equation\_wasserstein\_setting}), delivered together with the article's own complete original proof of it (source0.tex lines 438--481, one \texttt{proof} environment of 44 lines). No mathematics is written, repaired or re-derived: statement and proof are the article's own text line for line. Required context retained uncounted: lem:duality\_inequality (lines 345--350); assumption\_main (lines 391--398). Every definition, display and label of those lines is retained unchanged. Omitted from this package and not redistributed: 1--344, 351--390, 399--423, 436--437, 482--1430. Nothing inside a retained excerpt is omitted or altered: every retained body is delivered exactly as the article's lines are written, so no omission receipt of any kind accompanies this package. Citations: the retained excerpts carry no citation command at all, so no bibliography entry of the article is transcribed and no reference is redistributed. Reference adaptation: none was needed and none was made; every reference inside the delivered unit resolves inside it, so no unresolved reference or citation is produced. Printed slips kept literal: no printed word, symbol, hypothesis or number of the counted statement or of its proof was altered. Layout and class: the article's own class is \texttt{article} with packages including latexsym, amsfonts, amsmath, amssymb, epsfig, tabularx, amsthm, dsfont, mathrsfs, graphicx, hyperref, enumerate, color, enumitem; the wrapper is the lane's pinned \texttt{amsart} editorial wrapper at 10pt on A4, which restates the theorem-like environments the retained text needs and adds the article's own macro definitions that the retained text needs (\texttt{\textbackslash N}, \texttt{\textbackslash R}, \texttt{\textbackslash cA}, \texttt{\textbackslash cY}, \texttt{\textbackslash d}), with extra wrapper packages (bm, bbm, dsfont, xspace, cancel, mathtools). The delivered numbering is the unit's own. Assets: no retained span contains a figure environment, a graphics-inclusion command, a PDF-page inclusion, a table or a TikZ picture; no image file is copied into this package, the package declares no asset dependency and no image file is redistributed with it. Licence: version arXiv:2402.12122v2 of this article is distributed under the Creative Commons Attribution 4.0 International licence, as recorded in the arXiv licence metadata for this exact version and shown on the abstract page https://arxiv.org/abs/2402.12122v2. A full rights notice is delivered at licenses/arxiv-2402.12122-CC-BY-4.0.txt. Characters: every retained excerpt is pure ASCII, so the delivered engine, XeLaTeX with its default Latin Modern text font, reports no missing character.
\begin{lem}\label{lem:duality_inequality}
	For $\mathcal{W}$ w.r.t. distance-like function $d$ we have
	\[
	\mathcal{W}(\mu_1, \mu_2) \geq \sup_{\psi \in \text{Lip}_1} \left\vert  \int_\cA \psi(x) \mu_1(\d x) - \int_\cA \psi(y) \mu_2(\d y)  \right\vert.
	\]
\end{lem}

\begin{ass}\label{assumption_main}
	For each $\gamma \in \mathcal{I}$ we assume that $\pi_\gamma$ is the invariant distribution of $P_\gamma$ and that 
	\[
	\tau(P_\gamma) \leq M \qquad \text{and} \qquad
	\tau(P_\gamma^{k_0}) \leq \tau,
	\]
	with $M \in [1, \infty)$, $\tau \in [0,1)$ and $k_0 \in \N$ independent of $\gamma$.
\end{ass}

\begin{prop}\label{prop:poisson_equation_wasserstein_setting}
	Let $\gamma \in \mathcal{I}$, $f \colon \cY \to \R$ be $\pi_\gamma$ integrable and Lipschitz with constant $L \in (0, \infty)$ and let $E_\gamma (y) < \infty$ for any $y \in \cY$. 
	If Assumption~\ref{assumption_main} is true, then the function $u_\gamma\colon \cY \to \R$ given via 
	\[
	u_\gamma(y ) = \sum_{\ell=0}^\infty ( P_\gamma^\ell  f(y) - \pi_\gamma(f)) 
	\]   
	is well-defined and we have
	\begin{itemize}
		\item[i)] $\vert u_\gamma (y) \vert \leq \frac{k_0L M^{k_0} E_\gamma(y)}{1- \tau}$ for any $y \in \cY$,
		\item[ii)] $u_\gamma (y) - P_\gamma u(y) = f(y) - \pi_\gamma(f)$ for any $y \in \cY$.
	\end{itemize}
\end{prop}

\begin{proof}
	Let $\ell = tk_0 +r \in \N$ with $t \in \N_0$ and $0 \leq r < k_0$ as well as $y \in \cY$. 
	By Lemma \ref{lem:duality_inequality} and the properties of contraction coefficients   
	\begin{align*}
		\left\vert P_\gamma^\ell f(y) - \pi_\gamma(f)  \right\vert &\leq L \sup_{\psi \in \text{Lip}_1}\left\vert P^\ell \psi(y) - \pi_\gamma(\psi) \right\vert 
		\le L  \mathcal{W}(P_\gamma^{\ell}(y, \cdot), \pi_\gamma ) \\
		&\leq L \tau(P_\gamma^r) \tau(P_\gamma^{tk_0}) \mathcal{W}(\delta_y, \pi_\gamma)
		\leq L M^r \tau^{t} \mathcal{W}(\delta_y, \pi_\gamma) \\
		&\leq L M^{k_0} \tau^{t} E_\gamma(y) < \infty.
	\end{align*}
	Using this bound we obtain 
	\begin{align*}
		\vert u_\gamma(y )\vert  &\le \sum_{\ell=0}^\infty \vert P_\gamma^\ell f(y) - \pi_\gamma(f) \vert =
		\sum_{t=0}^\infty \sum_{j= tk_0 }^{(t+1) k_0-1 } \vert P_\gamma^{j} f(y) - \pi_\gamma(f) \vert\\
		&\leq \sum_{t=0}^\infty \sum_{j= tk_0 }^{(t+1)k_0 -1} L M^{k_0} \tau^{t} E_\gamma(y) 
		= \frac{k_0L M^{k_0} E_\gamma(y)}{1- \tau}.
	\end{align*}
	This shows that $u_\gamma$ is well-defined and also that i) holds. 
	
	
	It remains to show ii): 
	It is no restriction to assume that $\pi_\gamma (f) =0$ (otherwise consider $g = f- \pi_\gamma(f)$). 
	For any $n \in \N$ we set 
	\[
	\psi_n^+ = \sum_{\ell=0}^n P^\ell_\gamma f^+  \qquad\text{and}\qquad \psi_n^- = \sum_{\ell=0}^n P_\gamma^\ell f^-,
	\]
	where $f^+$ and $f^-$ are the positive and negative part of $f$. 
	For any $n \in \N$  the definition of $\psi_n^+$ implies
	\[
	P_\gamma \psi_n^+ = \psi_{n+1}^+ - f^+ \qquad\text{and}\qquad P_\gamma \psi_n^- = \psi_{n+1}^- - f^-.
	\]
	By monotone convergence for any $y_0 \in \cY$, 
	\[
	\lim_{n \to \infty} \left( P_\gamma \psi_n^+ (y_0) \right) = P_\gamma \left( \lim_{n \to \infty} \psi_n^+\right)  (y_0).
	\]
	By the same arguments, the same statement for $\psi_n^-$ (instead of $\psi_n^+$) is deduced.
	Clearly, $u_\gamma = \lim_{n \to \infty} (\psi_n^+ - \psi_n^-)$ and hence by linearity 
	\begin{align*}
		P_\gamma u_\gamma = P_\gamma \left( \lim_{n \to \infty} (\psi_n^+ - \psi_n^-) \right) = \lim_{n \to \infty} P_\gamma \psi_n^+  - \lim_{n \to \infty} P_\gamma \psi_n^- \\
		= \lim_{n \to \infty} \left(\psi_{n+1}^+ - \psi_{n+1}^-\right) - f
		= u_\gamma -f.
	\end{align*}
	This shows ii) and therefore completes the proof.
\end{proof}

\end{document}
